We measure the similarity of sixteen Philippine languages along three
independent dimensions: the characters in which they are written, the sounds in
which they are spoken, and the terrain across which they are distributed.
Working from Biblical translations collected by the authors, we compute
orthographic similarity as a Jaccard coefficient over trigram inventories and
phonetic similarity as a cosine over Double Metaphone frequency vectors, and we
interpret both against a custom linguistic map built from official atlases and
topographic data. The two quantitative measures agree closely in the ordering of
language relationships (Spearman \( \rho = 0.83 \)) while differing markedly in
scale: the phonetic band is roughly two-fifths as wide as the orthographic, a
disparity we show is not an artifact of corpus size. We read this as evidence
that orthography, the most recent and administratively tractable layer of a
language, registers colonial-era standardization at full amplitude, whereas
phonology preserves a more attenuated record of shared Austronesian inheritance.
Three findings recur across all three analyses: Chavacano is peripheral under
every measure, Romblomanon is central under every measure, and Northern Luzon
forms a distinct zone bounded by the Cordillera and Sierra Madre.
